\section*{NeurIPS Paper Checklist}

\begin{enumerate}

\item {\bf Claims}
    \item[] Question: Do the main claims made in the abstract and introduction accurately reflect the paper's contributions and scope?
    \item[] Answer: \answerYes{}
    \item[] Justification: The abstract reports multi-seed mean $\pm$ std results (seeds 42--46) that match Table~\ref{tab:full} exactly. The per-attack Table~\ref{tab:detection} clearly distinguishes multi-seed rows (ML only, ACB-Ref; seeds 42--46, 520 sessions, 8 attacks) from seed=42 supplementary rows (B7s, B8s, full-session LLM judge; labeled in caption). The contributions list in \S\ref{sec:intro} matches the content of \S\ref{sec:related}--\S\ref{sec:results}. The full-session LLM-as-judge results (61\% recall, 0\% hard FPR) are explicitly labeled as measured on ACB at seed=42. The total-FP equivalence finding is disclosed in the abstract and \S\ref{sec:results}.

\item {\bf Limitations}
    \item[] Question: Does the paper discuss the limitations of the work performed by the authors?
    \item[] Answer: \answerYes{}
    \item[] Justification: \S\ref{sec:gaps} (Open Gaps) covers: (1) B4 at 68\% requires an external agent-ID registry; (2) B5 at 70\% uses ML behavioral signal rather than a stateful settled-key registry; (3) B7s at 82.5\% requires merchant-domain knowledge beyond the current MCC allowlist; (4) EscRate (56\%) is a per-persona spend-ceiling design issue, not a calibration artifact---ACB-Ref hard FPR is exactly 0\%; (5) B9 excluded from this release; (6) B7s/B8s results are seed=42 only, not yet multi-seed. The full-session LLM judge (61\% recall) is described as an L0 ceiling subject to observation-surface limits, not as a deployable component.

\item {\bf Theory assumptions and proofs}
    \item[] Question: For each theoretical result, does the paper provide the full set of assumptions and a complete proof?
    \item[] Answer: \answerNA{}
    \item[] Justification: This is primarily a benchmark and systems paper. Definition~2.1 is a formal definition, not a theorem. No theorems are stated or proved. The paper's claims are empirical and definitional; any formal development is out of scope for this submission.

\item {\bf Experimental result reproducibility}
    \item[] Question: Does the paper fully disclose all information needed to reproduce the main experimental results?
    \item[] Answer: \answerYes{}
    \item[] Justification: The Reproducibility Statement gives the exact reproduction command for the core benchmark (\texttt{python -m benchmark.run\_experiment --seed 42 --no-sota}; CPU only, $<$30s/seed) and for the full-session LLM judge (\texttt{python -m benchmark.session\_fraud\_judge --all-attacks --seed 42 --n 20}; requires AWS Bedrock, $\approx$7 min, \$0.60). Markov parameters, amount distributions, perturbation $\sigma$, training seed, threshold formula ($\mu + 1.1\sigma$), and veto conditions are specified in \S\ref{sec:method}, \S\ref{sec:gen}, and the appendix. B7s/B8s attack generator is in \texttt{benchmark/synthetic.py}.

\item {\bf Open access to data and code}
    \item[] Question: Does the paper provide open access to data and code?
    \item[] Answer: \answerYes{}
    \item[] Justification: Generator, verifiers, and harness are released as open-source (anonymized URL in supplemental material; full URL in camera-ready). The benchmark is fully deterministic given \texttt{--seed}; no pre-collected dataset is required.

\item {\bf Experimental setting/details}
    \item[] Question: Does the paper specify all training and test details necessary to understand the results?
    \item[] Answer: \answerYes{}
    \item[] Justification: Data generation: persona-conditional Markov chains with parameters in \S\ref{sec:gen} and the appendix; perturbation $\sigma=0.05$. Train/test independence: ML trained on seed=7 ($N=500$), evaluated on seeds 42--46 ($N=520$ each). Threshold: $\mu+1.1\sigma$ of clean training scores. Veto: score$=0$ and uncertainty$>0.20$. All in \S\ref{sec:method} and \S\ref{sec:experiments}.

\item {\bf Experiment statistical significance}
    \item[] Question: Does the paper report error bars or appropriate information about statistical significance?
    \item[] Answer: \answerYes{}
    \item[] Justification: Table~\ref{tab:full} reports mean $\pm$ standard deviation across 5 independent seeds (42--46). The $\pm$ values are population standard deviations, not standard errors of the mean---stated explicitly in the table caption. All five seeds are fully independent (separate data generation, no shared state). Rates are bounded well away from 0 and 1 in the reported range, so symmetric error bars do not produce out-of-range values.

\item {\bf Experiments compute resources}
    \item[] Question: Does the paper provide sufficient information on compute resources needed to reproduce the experiments?
    \item[] Answer: \answerYes{}
    \item[] Justification: Stated in the Reproducibility Statement: CPU only (no GPU required), $<$30 seconds per seed, $<$3 minutes total for all 5 seeds. The Tier~3 LLM Verifier is invoked on $<$3\% of actions in the full pipeline but is bypassed by \texttt{--no-sota} for the main results; \texttt{--no-sota} results are the ones reported in all tables.

\item {\bf Code of ethics}
    \item[] Question: Does the research conform with the NeurIPS Code of Ethics?
    \item[] Answer: \answerYes{}
    \item[] Justification: This work promotes defensive safety in financial AI systems. All data is synthetically generated---no production transaction data or user data was used. No production payment credentials appear in the open-source release. The attack taxonomy is published at a conceptual level sufficient to motivate defenses, not at a working-exploit level.

\item {\bf Broader impacts}
    \item[] Question: Does the paper discuss both potential positive and negative societal impacts?
    \item[] Answer: \answerYes{}
    \item[] Justification: See the Broader Impact section following the Conclusion. Positive: protocol-layer safety tools may reduce financial fraud in agentic systems. Negative: the attack taxonomy, if studied adversarially, could inform attack design; mitigated by not publishing working exploit code and by releasing the detector alongside the taxonomy.

\item {\bf Safeguards}
    \item[] Question: Does the paper describe safeguards for responsible release of data or models with high risk for misuse?
    \item[] Answer: \answerYes{}
    \item[] Justification: ACB describes attack \emph{classes} at a conceptual and statistical level (Markov chains, amount distributions). No production payment API credentials, no merchant exploit scripts, and no working B5 settled-key replay against a live rail are included in the release. The released artifact is the detector (ACB-Ref), not the attack implementation. The benchmark generator produces synthetic sessions only.

\item {\bf Licenses for existing assets}
    \item[] Question: Are creators of assets used in the paper properly credited and license terms mentioned?
    \item[] Answer: \answerYes{}
    \item[] Justification: All assets cited are referenced, not redistributed --- no third-party code is reused in ACB. garak~\cite{garak} (Apache 2.0) and promptfoo~\cite{promptfoo} (MIT) are named as related tooling and are not run or bundled; Llama Guard~\cite{llamaguard} and FinHarness~\cite{finharness} are cited as prior work only. ISO~8583~\cite{iso8583} is a published standard (referenced, not reproduced); x402~\cite{x402}, its security analysis~\cite{x402attacks}, and the Stripe API reference~\cite{stripe} are cited as public specifications and documentation. The benchmarks we cross-validate against, AgentDojo~\cite{agentdojo} and $\tau$-bench~\cite{taubench}, are used under their published open-source licenses (MIT), with attack templates and task trajectories credited in \S\ref{sec:cross_bench}. All citations appear in the References section and were verified against their primary sources.

\item {\bf New assets}
    \item[] Question: Are new assets introduced in the paper well documented?
    \item[] Answer: \answerYes{}
    \item[] Justification: ACB (benchmark generator, verifiers, harness) is documented in \S\ref{sec:dataset} (methodology), the appendix (implementation details), and the Reproducibility Statement (exact run command). Released under MIT license. Anonymous supplemental URL at submission; full URL in camera-ready.

\item {\bf Crowdsourcing and research with human subjects}
    \item[] Question: Does the paper include full details of crowdsourcing or human subjects research?
    \item[] Answer: \answerNA{}
    \item[] Justification: All data is synthetically generated via parameterized Markov chains and log-normal amount distributions. No human subjects, no crowdsourcing, no real user data.

\item {\bf Institutional review board (IRB) approvals}
    \item[] Question: Does the paper describe potential risks to study participants and IRB approvals?
    \item[] Answer: \answerNA{}
    \item[] Justification: No human subjects research. Not applicable.

\item {\bf Declaration of LLM usage}
    \item[] Question: Does the paper describe the usage of LLMs if it is an important, original, or non-standard component of the core methods?
    \item[] Answer: \answerYes{}
    \item[] Justification: Two LLM uses are disclosed. (1) The Tier~3 LLM Verifier (\S\ref{sec:method}) uses \texttt{claude-haiku-4-5-20251001} as a classification component for high-uncertainty actions ($<3\%$ of actions in the full pipeline)---a method component described in \S\ref{sec:method}. (2) The full-session LLM-as-judge baseline uses \texttt{claude-sonnet-4-6} via AWS Bedrock to evaluate all 10 attack types ($n{=}20$/attack, seed=42, reported in Table~\ref{tab:detection})---an evaluation baseline, not a method component. LLMs are \emph{not} used in benchmark generation, L1 rules verification, or the ML behavioral model. LLMs were also used for writing assistance in preparing this manuscript, which does not require declaration per the NeurIPS handbook.

\end{enumerate}
